\documentclass[11pt]{article}
\usepackage[margin=1in]{geometry}
\usepackage{amsmath,amssymb,amsthm,booktabs}
\usepackage[hidelinks]{hyperref}
\setlength{\emergencystretch}{2em}
\newtheorem{theorem}{Theorem}
\newtheorem{proposition}[theorem]{Proposition}
\newtheorem{corollary}[theorem]{Corollary}
\newcommand{\mr}{\operatorname{mr}}
\newcommand{\rk}{\operatorname{rank}}
\newcommand{\F}{\mathbb F}
\newcommand{\Z}{\mathbb Z}
\title{Augmentation-module bounds from cover homology}
\author{Internal research working proof}
\date{16 November 2026}
\begin{document}
\maketitle

The user's proposed homology-gain test is valid. For an epimorphism
$G=\pi_1(\Sigma_2K)\twoheadrightarrow A_5$ and the inverse image
$H$ of a point stabilizer $A_4$, it gives
\[
 b_1(H;\F_p)-b_1(G;\F_p)\ge5,\quad p\ne5
 \quad\Longrightarrow\quad \mr(K)\ge4.
\]
It sharpens the earlier absolute Schreier threshold and does not
require a meridian-root lift to the knot group. All sixteen quotient
kernel classes on the retained sixteen-knot panel have now been
tested, including the four classes beyond the twelve old first maps.
Every cover has $H_1(H;\Z)=\Z/(3D)$ while $H_1(G;\Z)=\Z/D$,
with $\gcd(D,30)=1$. The gain is one at three and zero at every
other prime. None passes. A completeness argument below covers
\emph{every} connected index-five cover of these particular groups.
The full conjecture and all sixteen cases remain unsettled here.

\section{The general augmentation bound}

Let $H\le G$ have finite index $n>1$, and let $F$ be a field whose
characteristic does not divide $n$. The permutation module
$P=F[G/H]$ splits as
\[
 P=F\mathbf{1}\oplus V,\qquad
 V=\ker\bigl(\textstyle\sum:P\to F\bigr),\quad \dim V=n-1.
\]
Indeed the splitting of the sum map sends $a$ to $(a/n)\mathbf{1}$.
Transitivity makes $P^G=F\mathbf{1}$, and the sum of a constant
vector is $n$ times its coefficient, so $V^G=0$. No semisimplicity
assumption on the entire image group is needed.

\begin{theorem}\label{thm:gain}
Put $\Delta_F=b_1(H;F)-b_1(G;F)$. Then
\[
 \Delta_F=\dim_F H^1(G,V)\ge0,\qquad
 \rk G\ge1+\left\lceil\frac{\Delta_F}{n-1}\right\rceil.
\]
If $G=\pi_1(Y_\phi)$ is the simple degree-$r$ branched cover of a
knot associated to a transposition quotient, then
\[
 \mr(K)\ge r+\left\lceil\frac{\Delta_F}{n-1}\right\rceil.
\]
\end{theorem}
\begin{proof}
For finite index, the permutation module is the induced, equivalently
coinduced, trivial module of $H$. Shapiro's lemma gives
$H^1(G,P)\cong H^1(H,F)$. Apply cohomology to the displayed direct
sum to obtain the dimension identity. First cohomology with trivial
coefficients is dual to first homology over a field.

If $d$ arbitrary elements generate $G$, evaluation on those elements
injects the cocycle space $Z^1(G,V)$ into $V^d$: the cocycle rule
determines its values on every word. Thus $\dim Z^1\le d(n-1)$.
Principal cocycles have dimension $\dim V-\dim V^G=n-1$.
Subtracting gives $\Delta_F\le(d-1)(n-1)$, proving the rank bound.
Combine it with the companion cover inequality
$\mr(K)\ge r-1+\rk G$.
\end{proof}

For the five-point $A_5$ action, $n=5$ and $\dim V=4$. Thus
$\Delta_F\ge5$ forces $\rk G\ge3$. With $r=2$ this gives
$\mr(K)\ge4$; a four-seed bridge upper certificate gives equality.
Primes two and three are allowed, despite dividing $|A_5|$.
Only the index five needs to be invertible for this argument.

\paragraph{Comparison with Schreier.}
If $G$ were two-generated, Schreier gives $b_1(H;F)\le6$.
The new test improves this to
$b_1(H;F)\le b_1(G;F)+4$.
For $b_1(G;F)=0,1,2$, its respective sharp-four thresholds are
$b_1(H;F)\ge5,6,7$. The earlier absolute threshold seven is therefore
strictly weaker when the base Betti number is zero or one.

\section{Complete result on the retained panel}

The frozen selection is exactly the first sixteen names of the pinned
uncertainty panel, with every retained $A_5$ quotient-kernel class and
no adaptive target or map replacement. Their exact integral homology
and kernel counts are as follows; $\kappa$ is counted separately for
each knot.
\begin{center}
\begin{tabular}{lrr|lrr}
\toprule
Knot & $D$ & $\kappa$ & Knot & $D$ & $\kappa$\\
\midrule
$13a129$&391&1&$13a233$&403&0\\
$13a135$&401&2&$13a234$&389&1\\
$13a136$&391&1&$13a239$&361&0\\
$13a137$&401&2&$13a242$&359&1\\
$13a194$&391&1&$13a243$&337&0\\
$13a196$&401&2&$13a244$&311&1\\
$13a197$&391&1&$13a253$&337&0\\
$13a198$&401&2&$13a265$&311&1\\
\bottomrule
\end{tabular}
\end{center}
The base homology is $\Z/D$. For every one of the sixteen classes
on twelve knots, the index-five subgroup has homology $\Z/(3D)$.
In particular the four extra maps on $13a135,13a137,13a196,13a198$
do not change the outcome. Since every $D$ is coprime to thirty,
\[
\begin{array}{c|ccc}
 p&3&p\mid D&\text{all other primes}\\\hline
 b_1(G;\F_p)&0&1&0\\
 b_1(H;\F_p)&1&1&0\\
 \Delta_{\F_p}&1&0&0.
\end{array}
\]
This is an all-prime conclusion from integral Smith forms, not a
finite prime sweep. The four knots with $\kappa=0$ have no $A_5$
map from which to form the proposed cover.

\begin{proposition}\label{prop:allfive}
For these sixteen base groups, every connected index-five cover is
represented by the tested classes. Consequently no such cover
satisfies the user's Betti-gain criterion over any prime field with
$p\ne5$.
\end{proposition}
\begin{proof}
An index-five subgroup gives a transitive coset action $G\to S_5$.
Its sign is trivial because the abelianization $\Z/D$ has odd order,
so its image lies in $A_5$. Every proper nontrivial subgroup of
$A_5$ has a nontrivial abelian quotient of order divisible by two,
three or five. This elementary finite subgroup fact is checked on
the complete 59-subgroup lattice in the retained arithmetic receipt.
Such an image would give $G$ a nontrivial abelian quotient at one
of these primes, contrary to $\gcd(D,30)=1$. Transitivity rules out
the trivial image. The image is therefore all of $A_5$.

All five subgroups $A_4\le A_5$ are conjugate. Postcomposing an
epimorphism with an automorphism of $A_5$ changes its point-stabilizer
preimage only by conjugacy in $G$. Hence one representative per
quotient kernel suffices for the unbased connected covers. The
earlier complete enumeration gave exactly the classes tested here.
\end{proof}

The cyclic abelianizations do not imply that the subgroups themselves
are cyclic. The proposition excludes this particular homological
certificate in degree five, not other cover degrees, nonabelian
subgroup information or new representations.

\section{Replay, controls and comparison with old exclusions}

The producer builds all 70 relation rows of the five-sheeted cover
of the normalized thirteen-generator core presentation, collapses a
four-edge spanning tree, and computes Smith certificates on the
remaining 61 columns. It checkpoints after every cover under a frozen
sixteen-cover cap. Previous first-witness evidence is retained.

The verifier independently builds the integer matrix by evaluated
Fox prefixes and replays every elementary Smith operation. It also
constructs a separate cocycle matrix directly on the basis
$e_i-e_5$, $1\le i\le4$, of the augmentation module. For each
relator, the coefficient of a positive letter is its prefix action;
for an inverse letter it is minus the action after that inverse.
Every prime dividing a retained torsion factor and one generic
prime are checked, with five excluded. The cocycle dimension minus
four agrees with the homology gain in each case. The integral Smith
data handle all other primes. The PD parser and Smith replay are
shared; the three matrix constructions are not described as wholly
independent software.

Three controls show the threshold and its boundary:
\begin{center}
\begin{tabular}{lccc}
\toprule
Base group & Index-five $H_1(H;\Z)$ & $b_1(G;\F_3)$ & $\Delta_{\F_3}$\\
\midrule
$F_2$ & $\Z^6$ &2&4\\
$C_3*C_3$ & $(\Z/3)^4$ &2&2\\
$C_3*C_3*C_3$ & $\Z^2\oplus(\Z/3)^6$ &3&5\\
\bottomrule
\end{tabular}
\end{center}
The third is the double-cover group of a connected sum of three
trefoils, an already-known sharp four-bridge control. The first
attains the upper bound four for a two-generated group.

At characteristic five, the same $F_2$ action gives a necessary
negative control: the augmentation module contains the constant
vector, has a one-dimensional invariant space and three-dimensional
principal cocycles. Therefore $\dim H^1(F_2,V)=8-3=5$, although
$b_1(H)-b_1(F_2)=6-2=4$. The claimed splitting identity fails there;
the prime restriction cannot simply be dropped.

The user's distinction from the old exclusions is correct. Those
concern specific knot-group representations with prescribed meridian
defect and root geometry. Here the group is the \emph{double-cover}
group, and the module is its four-dimensional permutation
augmentation module. No knot-group extension is imposed. The old
$H_3$ root rigidity and the statement comparing full versus
restricted cocycles on the \emph{same} representation do not exclude
this new test. The direct computation above supplies its negative
result. The receipts are \path{augmentation_cover_verification.json}
and the unchanged earlier \path{joint_a5_verification.json}.

\section{Whitehead dependency and a peripheral obstruction}

A fresh retrieval of the external project's open-list blob agrees
with the retained 293-name file. All sixteen targets are in it,
and none is among the separate 243 certificates. This supports the
user's status correction within that release; it is not an assertion
that all subsequent literature is exhausted. The unrefereed
certificates retain their earlier algebra, name and bridge replay.

The new preprint of Gao--Gu--Li--Wang supplies a different mechanism
for iterated untwisted Whitehead doubles of two-bridge knots
\cite{gglw}. Its peripheral curve condition uses orders
$\nu(\mu)=-p$, $\nu(\eta)=q$ with $p>0$ and either $q>p$ or $q<0$,
where $\eta=\lambda^2$. The points must come from actual
augmentations. This is not an arbitrary choice of peripheral
eigenvalues or a closure of the sixteen targets. The one-step
construction in Section 2.3 requires $\eta\ne1$.

\begin{proposition}\label{prop:orthogonal}
Over a field of characteristic different from two, suppose a knot
representation preserves a nondegenerate symmetric form, and its
meridian is diagonalizable with nonidentity defect of rank one.
On its distinguished eigenline, the longitude eigenvalue $\lambda$
satisfies $\lambda^2=1$. Thus this representation lies on the
excluded $\eta=1$ locus of that Whitehead construction.
\end{proposition}
\begin{proof}
Write $V=L\oplus F$ for the nonidentity meridian eigenline and its
fixed hyperplane. Invariance of the form makes $L$ orthogonal to
$F$. Nondegeneracy then makes its restriction to $L$ nonzero;
otherwise $L$ would be orthogonal to all of $V$. Since the longitude
commutes with the meridian, it preserves $L$. Preservation of the
form on this line gives $\lambda^2=1$.
\end{proof}

Consequently a fixed orthogonal reflection base cannot simply seed
the peripheral construction. Any useful deformation must leave this
orthogonal constraint and supply actual representation data.
Transvections are a separate excluded input: their meridian eigenvalue
is one, whereas the cited construction requires a diagonalizable
nonidentity eigenvalue. These observations do not prohibit
deformations to other representations; no such deformation has been
constructed here.

\paragraph{Original Dutra lemma.}
The reported gap can be checked directly in \cite[p.~7]{dutra}.
In the stated free group on $x_1,x_2$, its embeddings give
\[
 \alpha(y_1y_2)=\omega(y_1y_2)
   =x_2x_1^{-1}x_2^{-1}x_1\ne1.
\]
This freely reduced four-letter word belongs to both $U_\alpha$
and $U_\omega$, contradicting the trivial intersection claimed in
Lemma 3.1 at conjugator one. Page twelve explicitly uses that lemma
in the proof of Lemma 3.8(2). Both complete pages were visually
checked and the free reduction is retained in the dependency audit.
This disproves the lemma as stated, not the main Whitehead-double
MRC conclusion. No proved record in this programme uses Dutra's
2022 Theorem 1.1 as a dependency. Its source row remains flagged;
the published 2017 graph-manifold/braid-pattern result is a distinct
earlier theorem. The new preprint explicitly avoids these subgroup
lemmas. The complete 44-page proof remains unaudited here.

Uniform arguments in this note are working proofs without formal
verification or external specialist review. The finite all-index-five
exclusion and peripheral obstruction are progress boundaries, not
a full MRC proof or a new resolved knot. Further work requires
another degree or module with a structural rank-gain reason, or an
actual peripheral representation curve outside the stated constraints.

\clearpage
\begin{thebibliography}{9}
\bibitem{gglw} Gao--Gu--Li--Wang, Meridional ranks of iterated
Whitehead doubles, preprint, 11 November 2026. Introduction footnote;
Sections 2.3--2.4; Definition 4.1, Theorem 4.2 and Corollary 4.5.
\url{https://arxiv.org/html/2611.15673v1}.
\bibitem{dutra} Dutra, Meridional rank of Whitehead doubles,
preprint v1, 26 December 2022, pp.~7 and 12.
\url{https://arxiv.org/abs/2212.13081v1}.
\bibitem{core} Silver--Traldi--Williams, Core groups (2024),
Section 7. \url{https://arxiv.org/abs/2311.02048v2}.
\end{thebibliography}
\end{document}
